%!TEX root = main.tex
%=========================================================

\section{Related Work}
\label{sec:related}

We review related work on broadcast in MANETs, adaptation mechanisms in wireless and mobile networks, and studies of mobility in these systems.

A survey by Ruiz and Bouvry~\cite{SurveyRuiz} provides a comprehensive coverage of broadcast in MANETs.
It provides a taxonomy of controlled flooding and overlay-based protocols, also covered respectively in surveys by Reina \emph{et~al.}~\cite{Reina2015} and Santi~\cite{Santi:2005:TCW:1089733.1089736}.
Approaches are subdivided based on the use of fixed- or variable-range wireless communication.
We considered in this paper only the former class of protocols~\cite{BroadcastStormProblem,ProbFlood,stojmenovic-cds,Nikolov2015,mpr,mpr-improvement,overlays_for_manets} and leave the adaptation of variable-range broadcast~\cite{woon2011variable,park2013efficient} to future work.

\cite{BroadcastStormProblem}
\cite{2001:AAR:876878.879323}: the one he uses, controlled flooding with timer and counters. Also proposes two others: using timers based on surface and others using signal noise threshold (both we assume we don't have)
\er{we need two more papers of recent work on flooding protocols}

overlay: also considered a \emph{connected dominated set} (CDS). Was pioneered by~\cite{stojmenovic-cds}, who is using the calculation of 2-hop neighbors, and calculating differences in covered nbors sets.
Several other protocols proposed to construct CDS.
\cite{mpr}: the one we use
\cite{mpr-improvement}: improvements, we do not use. Detail?
\cite{overlays_for_manets}: uses geolocation, based on the knowledge of the transmission areas of each node and neighbor sets
\cite{Nikolov2015}: keeps history of past nbors, determines to switch to relay based on coverage of nodes in the recent past

\bigskip

broadcast algorithms in manets
\begin{itemize}
	\item Surveys:
	\begin{itemize}
		\item : 2015 survey with the taxonomy
		\item \cite{Reina2015}: 2015 survey on probabilistic broadcast approaches 
		\item \cite{Santi:2005:TCW:1089733.1089736}: 2005 survey on topology control in manets
	\end{itemize}
	\item Controlled flooding:
	\begin{itemize}
		\item \cite{BroadcastStormProblem}: 3 controlled flooding approaches to deal with broadcast storms
		
		% \item \cite{ProbFlood}: 2003 broadcast approach to privilege that nodes located at the border of transmissions range act as relays
	\end{itemize}
	\item Overlay-based
	\begin{itemize}
		\item \cite{stojmenovic-cds} 1st paper in determining connected dominating sets (CDS)
		\item \cite{Nikolov2015} 2015 a more recent approach to build CDS
		\item \cite{mpr} another approach to create CDS
		\item \cite{mpr-improvement}: 2014 recent improvement of MPR (multi-point relaying technique) to build backbones
		\item \cite{overlays_for_manets} 2015 recent study about how to maintain P2P overlays in MANETs for routing
	\end{itemize}

	%		\item \cite{clustering-approach}: 2009 approach for clusters within MANETs that uses location and energy metrics \er{for clusters? from the title it seems it uses broadcast as a tool, so you could rather use it in the intro to motivate the uses of broadcast}

\end{itemize}

studies of broadcast performance
\begin{itemize}
	\item \cite{DensityMobDrivenEval} our evaluation study of broadcast over different scenarios of mobility and density
% \raziel{not required because Ruiz survey is more recent than this study} \er{strongly disagree, this is the only paper with the exception of ours that performs a performance evaluation, we must cite it at least in the introduction.}
	\item evaluation of broadcast protocols~\cite{Williams2002}, there is also a first classification of protocols
\end{itemize}

adaptation mechanisms in wireless systems, manets, sensor networks
\begin{itemize}
	\item \cite{adaptive-multihop}
	\item \cite{adaptive_group_com_iscc02} is basically one adaptation policy where the switching criteria is based on the relative velocity of nodes ; nodes switch between a list of controlled flooding approaches
	\item in this work \cite{Bellavista2013} overlays are build on the top of wireless sensor networks (WSN) to speed up data dissemination. Such data is classified with a certain priority; interoperability issues between WSN and MANETs are tackled but no between dissemination protocols
	\item 2003 \cite{2001:AAR:876878.879323} this work (aprox 700 citaitons) follows the paper about the broadcast storm problem to point out the need of adapting the value of thresholds shown in \cite{BroadcastStormProblem}
	\item \cite{hypergossiping} : adaptive approach based on gossiping
	\item \cite{prob-based-adaptive-tech} thresholds are changed on the fly without using control messages (seems promising, still have to read carefully)
	\item \cite{whitbeck2010hymad}: adaptation for routing in dense and dynamic networks
\end{itemize}

mobility traces and models
\begin{itemize} 
	\item \cite{Hess:2015:DHM:2856149.2840722}: a recent survey
	\item \cite{Aschenbruck2009}: mobility model for disaster area scenarios
	\item \cite{rhee2011levy}: Levy walk nature, classical paper
\end{itemize}

emergent overlays elsewhere
\begin{itemize}
	\item This paper~\cite{galuba2007generic} uses the term in another context, an unstructured p2p network where a DHT or another structure is bootstrapped on demand
\end{itemize}